\documentclass[11pt]{article}
\usepackage[margin=1in]{geometry}
\usepackage{amsmath}
\usepackage{booktabs}

\title{Step 2 Validation Protocol: Separating Constant SOC, Dimerisation, and Local Spin-Correlation Controls}
\author{}
\date{}

\begin{document}
\maketitle

\section{Purpose of the test}
The purpose of this step is to determine whether the minimal bright/dark
spectroscopy model can distinguish three phenomenological mechanisms for the
effective bright-dark mixing:
\begin{align}
V_{BD}^{\mathrm{eff}} &= V_0,\\
V_{BD}^{\mathrm{eff}} &= V_0 + \lambda_\delta \delta,\\
V_{BD}^{\mathrm{eff}} &= V_0 + \lambda_C C_1.
\end{align}
The test is not a microscopic calculation of CuGeO3.  It is a controlled
spectroscopic check built on the Step 1 model.

\section{Physical model}
The Hilbert space contains a ground state, a dark excited state, and a bright
excited state.  The bright state carries the optical dipole by default.  The
dark state becomes visible only through bright-dark mixing.  The effective
mixing is
\begin{equation}
V_{BD}^{\mathrm{eff}} =
V_0 + \lambda_\delta \delta + \lambda_C C_1.
\end{equation}
Here $\delta$ is a spin-Peierls/lattice dimerisation order parameter, while
$C_1=\langle S_i\cdot S_{i+1}\rangle$ is a local spin correlation that can be
nonzero even in the uniform chain.

\section{Expected results}
If $V_0$ alone is active, the cross-peak signal should reproduce the qualitative
Step 1 bright/dark result.  If only $\lambda_\delta\delta$ is varied, changes in
the cross peaks should track the dimerisation amplitude.  If only
$\lambda_C C_1$ is varied, cross-peak changes can remain present even when
$\delta=0$, which would distinguish spin-correlation control from a purely
spin-Peierls order parameter effect.

\section{Numerical procedure}
Run one scenario at a time by changing \texttt{active\_case\_key} in the
notebook or script.  The third-order rephasing 1Q spectrum is generated with
the same NQ protocol used in Step 1.  Each run writes a spectrum PDF, an NPZ
file containing the complex response arrays, and text/JSON summaries.

\section{Acceptance criteria}
This step is qualitative and comparative.  The model is considered useful for
Step 2 if the three controls produce distinguishable changes in at least one of:
cross-peak area, cross/diagonal ratio, peak position, or pathway-resolved
amplitude.  Absolute convergence in frequency-grid resolution is not required
for this first separation test.

\section{Requested figures and data}
For each scenario, save the rephasing real spectrum, the pathway-resolved
response arrays, the total rephasing component, and the run parameters.  Later
Phase 2 analysis should compare cross-peak finite-window integrals and
cross/diagonal ratios across the three scenarios.

\section{Interpretation of possible failures}
If all three controls produce nearly identical spectra at matched
$V_{BD}^{\mathrm{eff}}$, then the current observable is sensitive only to the
total effective mixing, not to its microscopic origin.  If the dimerisation and
spin-correlation controls differ only through amplitude, additional observables
such as temperature/field trends, unrephasing pathways, or dynamic phonon
coupling may be required.

\section{Numerical results}
The three requested scenarios were generated at $\eta=0.001$ with a
$100\times100$ frequency grid.  The finite-window half-width used for the
aggregate metrics is $0.025$ eV.  The reported cross area is the sum of the
two cross windows, and the ratio is the cross-window area divided by the sum
of the two diagonal-window areas.

\begin{center}
\begin{tabular}{lrrrr}
\toprule
Scenario & $V_{BD}^{\mathrm{eff}}$ & $\delta$ & $C_1$ &
$A_{\mathrm{cross}}/A_{\mathrm{diag}}$ \\
\midrule
Constant SOC & 0.0100 & 0 & 0 & 0.120590 \\
Dimerisation control & 0.0200 & 0.0500 & 0 & 0.105788 \\
Spin-correlation control & 0.0200 & 0 & -0.443 & 0.105788 \\
\bottomrule
\end{tabular}
\end{center}

The summed cross-window finite-area values are $10.210040$ for constant SOC
and $13.682378$ for both enhanced-mixing controls.  The direct spectrum
difference between the dimerisation and spin-correlation controls at matched
$V_{BD}^{\mathrm{eff}}=0.02$ has relative maximum difference
$2.62\times10^{-17}$.

\section{Validation outcome}
The first criterion passes: the model remains sensitive to the effective
bright-dark mixing, since the constant-SOC case and the enhanced-mixing cases
produce distinguishable cross-window amplitudes.

The microscopic-origin criterion is inconclusive in this scalar model.  The
dimerisation and spin-correlation controls are numerically identical at matched
$V_{BD}^{\mathrm{eff}}$, because $\delta$ and $C_1$ enter the Hamiltonian only
through the same scalar effective mixing.  Therefore this Step 2 model can
separate weak and enhanced bright-dark mixing, but it cannot by itself assign
the enhancement to spin-Peierls dimerisation rather than to local spin
correlations.

\section{Conclusions}
This step gives a useful negative control.  The nonlinear response is clearly
sensitive to the magnitude of the effective bright-dark mixing: increasing
$V_{BD}^{\mathrm{eff}}$ from $0.01$ to $0.02$ changes the cross-window signal
and the full complex spectrum.  This confirms that the Step 1 observable remains
usable when the mixing is promoted from a fixed SOC parameter to an effective
control parameter.

However, the present implementation cannot distinguish the microscopic origin
of the enhancement.  In the Hamiltonian used here, $\delta$ and $C_1$ only enter
through the single number $V_{BD}^{\mathrm{eff}}$.  Consequently, a
dimerisation-driven enhancement and a local-spin-correlation-driven enhancement
are exactly equivalent when they produce the same $V_{BD}^{\mathrm{eff}}$.  The
matched comparison gives a relative maximum spectral difference of
$2.62\times10^{-17}$, which is numerical zero for the present calculation.

The physical conclusion is therefore limited but important: a temperature- or
field-dependent cross peak cannot be assigned to SOC, dimerisation, phonons, or
spin correlations from this scalar-mixing model alone.  The model can detect an
enhanced bright-dark mixing channel, but it cannot identify whether the
enhancement is caused by spin-Peierls lattice order or by correlations already
present in the uniform spin chain.

\section{Proposed next steps}
\subsection{Step 3A: add a dynamic dimerisation/phonon coordinate}
The most direct next step is to give the dimerisation channel a dynamical
signature rather than only a scalar amplitude.  A minimal implementation would
keep the bright/dark sector but add an explicit phonon coordinate $Q$:
\begin{equation}
H_{\mathrm{mix}} =
\left(V_0 + \lambda_\delta \delta\right)
\left(|D\rangle\langle B| + |B\rangle\langle D|\right)
+ g_Q Q
\left(|D\rangle\langle B| + |B\rangle\langle D|\right).
\end{equation}
This should create phonon sideband structure, modified pathway phases, or
additional dephasing/lineshape changes that a static $C_1$ control does not
produce.  This is the cleanest route for separating lattice-driven
spin-Peierls physics from a purely scalar SOC enhancement.

\subsection{Step 3B: add independent temperature and field trends}
The second control should be a trend-level comparison:
\begin{align}
\delta(T,B) &\rightarrow 0 \quad \text{above the spin-Peierls transition},\\
C_1(T,B) &\neq 0 \quad \text{in both uniform and dimerised regimes}.
\end{align}
If a spectral feature follows $\delta(T,B)$ and vanishes above the transition,
it is compatible with spin-Peierls/lattice control.  If it persists when
$\delta=0$ but $C_1$ remains finite, it should not be attributed uniquely to
the dimerised phase.

\subsection{Step 3C: introduce a k- or bond-sector distinction}
The third step is to make the dimerised and uniform spin-chain controls act on
different sectors.  For example, a dimerised phase can fold the Brillouin zone
or mix $k$ and $k+\pi$, while a uniform-chain correlation control should preserve
the original translation symmetry.  This would test whether rephasing and
unrephasing spectra are sensitive to a genuine spin-Peierls sector structure,
not only to a scalar change in $V_{BD}^{\mathrm{eff}}$.

\subsection{Step 3D: compare more observables than cross-peak area}
The Step 2 area metric is useful but not sufficient.  Future runs should compare
at least:
\begin{itemize}
\item rephasing versus unrephasing totals;
\item real, imaginary, and absolute spectra;
\item cross/diagonal ratios;
\item pathway-resolved amplitudes;
\item peak positions and linewidths;
\item possible phonon sidebands or asymmetric lineshapes.
\end{itemize}
The goal is to find an observable that changes under dimerisation/phonon
dynamics but does not change under a purely static $C_1$ scalar control.

\subsection{Recommended order}
The recommended next model is Step 3A: keep the same bright/dark model and add
one explicit dimerisation/phonon coordinate.  This is the smallest change that
can break the degeneracy found in Step 2.  After that, Step 3B should test
whether the resulting spectral signatures follow realistic temperature and
field trends for $\delta(T,B)$ and $C_1(T,B)$.

\end{document}
